% Part: many-valued-logic
% Chapter: three-valued-logics
% Section: lukasiewicz

\documentclass[../../../include/open-logic-section]{subfiles}

\begin{document}

\olfileid{mvl}{inf}{luk}

\olsection{\L ukasiewicz 逻辑}

\begin{editorial}
  本节是关于无限值 \L ukasiewicz 逻辑的一个简短“存根”。
\end{editorial}

\begin{defn}\ollabel{def:lukasiewicz}
无限值 \L ukasiewicz 逻辑~$\LogLuk[\infty]$ 使用如下矩阵定义：
\begin{enumerate}
  \item 标准命题语言 $\Lang L_0$，带有联结词
  $\lnot$、$\land$、$\lor$、$\lif$。
  \item 真值集为 $V_\infty$。
  \item $1$ 是唯一指定值，即 $V^+ = \{1\}$。
  \item 真值函数由以下等式给出：
  \begin{align*}
    \tf{\lnot}[\LogLuk](x) & = 1 - x\\
    \tf{\land}[\LogLuk](x,y) & = \min(x,y)\\
    \tf{\lor}[\LogLuk](x,y) & = \max(x,y)\\
    \tf{\lif}[\LogLuk](x,y) & = \min(1,1-(x-y)) = \begin{cases}
      1 & \text{若 } x \le y\\
      1-(x-y) & \text{否则。}
    \end{cases}
    \end{align*}
\end{enumerate}
$m$ 值 \L ukasiewicz 逻辑的定义与此相同，只是 $V = V_m$。
\end{defn}

\begin{prop}
  由 \olref[thr][luk]{def:lukasiewicz} 定义的逻辑 $\LogLuk[3]$ 与
  由 \olref{def:lukasiewicz} 定义的逻辑 $\LogLuk[3]$ 相同。
\end{prop}

\begin{proof}
  通过比较 \olref[thr][luk]{def:lukasiewicz} 中给出的联结词真值表与
  \olref{def:lukasiewicz} 中等式所确定的真值表即可看出：
  \begin{center}
    \begin{tabular}{c|c}
      $\tf{\lnot}$ & \\
      \hline
      $1$ & $0$ \\
      $1/2$ & $1/2$ \\
      $0$ & $1$
    \end{tabular}
    \quad
    \begin{tabular}{c|ccc}
      $\tf{\land}[\LogLuk[3]]$ & $1$ & $1/2$ & $0$ \\
      \hline
      $1$ & $1$ & $1/2$ & $0$ \\
      $1/2$ & $1/2$ & $1/2$ & $0$\\
      $0$ & $0$ & $0$ & $0$
    \end{tabular}
    \\[2ex]
    \begin{tabular}{c|ccc}
      $\tf{\lor}[\LogLuk[3]]$ & $1$ & $1/2$ & $0$ \\
      \hline
      $1$ & $1$ & $1$ & $1$ \\
      $1/2$ & $1$ & $1/2$ & $1/2$ \\
      $0$ & $1$ & $1/2$ & $0$
    \end{tabular}
    \quad
    \begin{tabular}{c|ccc}
      $\tf{\lif}[\LogLuk[3]]$ & $1$ & $1/2$ & $0$ \\
      \hline
      $1$ & $1$ & $1/2$ & $0$ \\
      $1/2$ & $1$ & $1$ & $1/2$  \\
      $0$ & $1$ & $1$ & $1$
    \end{tabular}
  \end{center}
\end{proof}

\begin{prop}\ollabel{prop:luk-infty-m}
  如果 $\Gamma \Entails[\LogLuk[\infty]] !B$，那么对所有 $m \ge 2$，
  都有 $\Gamma \Entails[\LogLuk[m]] !B$。
\end{prop}

\begin{proof}
  练习。
\end{proof}

\begin{prob}
  证明 \olref[mvl][inf][luk]{prop:luk-infty-m}。
\end{prob}

事实上，其逆也成立。

无限值 \L ukasiewicz 逻辑是最主流的模糊逻辑。
在模糊逻辑文献中，条件句常被定义为 $\lnot !A \lor !B$。由此得到的结果将是无限值强 克林逻辑。

\begin{prob}
  证明 $(p \lif q) \lor (q \lif p)$ 是 $\LogLuk[\infty]$ 的一个重言式。
\end{prob}

\end{document}